\subsection{Convergencia de CV\% de EDP -- la laguna declarada en la sección anterior}

La sección anterior dejó explícitamente sin resolver una pregunta:
\emph{"no hay evidencia todavía de que tres repeticiones sean suficientes
para estimar el EDP con la misma confianza"} que el tiempo de ejecución.
Este apartado la cierra, usando el dataset final de CPU ya corregido
(\texttt{arc174}, tras el fix de \texttt{repetition\_index} del
2026-08-25) en vez del barrido dedicado: 9 kernels $\times$ 6 niveles de
frecuencia $\times$ 10 repeticiones reales, EDP $=$ energía $\times$
tiempo por repetición, mismo estimador (CV\% con desviación poblacional)
y mismo protocolo (orden fijo $n=3,6,10$ más la distribución combinatoria
completa sobre $\binom{10}{3}=120$ tríos) que la Sección anterior.

\textbf{La convergencia de EDP es peor que la de tiempo, no igual.} De las
54 celdas (kernel, nivel), 48 ya son estables (CV\%$<$2\,\%) desde $n=3$ y
se mantienen así hasta $n=10$ -- estas replican el patrón ya visto para
tiempo. Pero dos celdas revelan varianza real que ni siquiera $n=10$
termina de resolver:

\begin{itemize}
  \item \textbf{\texttt{npb\_ft@REF}}: CV\% =1{,}57\,\% a $n=6$, pero
    \textbf{sube a 6{,}46\,\% a $n=10$} -- el CV\% no converge de forma
    monótona, sino que una repetición más tardía revela dispersión que
    las primeras seis no mostraban. El mismo patrón se repite en
    \texttt{npb\_ft@F0} y \texttt{npb\_ft@F1}.
  \item \textbf{\texttt{rodinia\_lavamd\_omp@F1}}: CV\% =1{,}26\,\% a
    $n=3$ (orden fijo), salta a 8{,}08\,\% a $n=6$ y se estabiliza en
    6{,}63\,\% a $n=10$ -- con solo tres repeticiones esta dispersión
    real nunca se habría observado, el mismo modo de falla ya documentado
    para \texttt{rodinia\_lud} en la sección de tiempo.
\end{itemize}

El peor caso sobre la distribución combinatoria completa a $n=3$ (no solo
el orden fijo) es \texttt{rodinia\_lavamd\_omp@F1} con un percentil 95 de
10{,}06\,\% -- es decir, un trío con mala suerte de tres repeticiones
detectaría una dispersión de EDP de dos dígitos, mientras que el trío que
efectivamente se habría ejecutado primero (orden fijo) la habría ocultado
casi por completo.

\begin{table}[H]
\centering
\caption{Cobertura acumulada: número de las 54 celdas (kernel, nivel) con CV\% de EDP por debajo de un umbral, según $n$.}
\label{tab:edp-convergence-coverage}
\begin{tabular}{@{}lrrr@{}}
\toprule
\textbf{Umbral} & \textbf{n=3 (fijo)} & \textbf{n=6 (fijo)} & \textbf{n=10} \\
\midrule
CV\% $<$ 2\,\%  & 48/54 & 48/54 & 45/54 \\
CV\% $<$ 5\,\%  & 52/54 & 50/54 & 50/54 \\
CV\% $<$ 10\,\% & 54/54 & 54/54 & 54/54 \\
\bottomrule
\end{tabular}
\end{table}

Que la cobertura a umbral 2\,\% \emph{empeore} de $n=3$ (48/54) a $n=10$
(45/54) no es una regresión del instrumento: es la misma advertencia de
la sección anterior, confirmada aquí con más fuerza. Una muestra pequeña
subestima sistemáticamente la varianza real; agregar repeticiones no
"introduce" dispersión, la \emph{revela}. El costo de esa revelación es
que ninguna $n$ fija y pequeña puede prometerse suficiente para todos los
kernels: \texttt{npb\_ft} todavía no muestra señales claras de haber
convergido en $n=10$, el máximo disponible en este dataset.

\textbf{Conclusión operativa, distinta de "usar más repeticiones para
todos".} La mayoría del catálogo (48 de 54 celdas, el 89\,\%) no necesita
diez repeticiones para estimar su EDP con confianza -- $n=6$ conserva
exactamente la misma cobertura a umbral 2\,\% que $n=3$ y $n=10$ (48/54),
por un 40\,\% menos de costo que diez. El problema no está distribuido
uniformemente: está concentrado en dos kernels (\texttt{npb\_ft},
\texttt{rodinia\_lavamd\_omp} en niveles intermedios), que exhiben
varianza de EDP genuina y no completamente caracterizada ni siquiera con
el máximo $n$ disponible. Forzar $n=10$ para las 54 celdas para blindar
dos de ellas es más caro que necesario; recortar a $n=3$ para todas
ocultaría exactamente la dispersión que estas dos ya demostraron tener.

\textbf{Recomendación para la ampliación del catálogo CPU (kernels
sobrevivientes del tamizaje RAJAPerf, ver Anexo N.2 del plan maestro de
Fase 2):} usar $n=6$ como protocolo por defecto -- evidencia-basado, no
heredado sin más de Fase 1 ni recortado a ciegas por presión de tiempo -- y
tratar cualquier kernel nuevo cuyo CV\% de tiempo u OI durante el
tamizaje ya sugiera comportamiento inestable (I/O significativo,
dependencia de datos, saturación temprana) como candidato a repeticiones
adicionales dirigidas, en vez de inflar el protocolo completo. Para
\texttt{npb\_ft} y \texttt{rodinia\_lavamd\_omp} específicamente, que ya
forman parte del catálogo confirmado, esta varianza no resuelta se declara
como limitación conocida del dataset, no se oculta ni se promete resuelta
con más repeticiones sin evidencia de que eso funcione.
